

\illustration{.3}{CuttySarkRomNum}{\href{https://en.wikipedia.org/wiki/Roman_numerals\#/media/File:CuttySarkRomNum.jpg}{Roman numerals on stern of the ship Cutty Sark showing draught in feet. The numbers range from 13 to 22, from bottom to top.} (CC BY-SA 3.0)}

The Roman numerals are a numeral system that originated in ancient Rome and was widely used throughout Europe well into
the Late Middle Ages. It differs from the Arabic system that we mostly use today in that numbers are written with combinations
of letters from the Latin alphabet, where each letter is assigned a fixed integer value.  Over the years many different letters were used,
which was a cause of frequent confusion.  But ``modern'' style uses only these seven:

\begin{tabular}{ccccccc}
\textbf{I} & \textbf{V} & \textbf{X} & \textbf{L} & \textbf{C} & \textbf{D} & \textbf{M} \\
1 & 5 & 10 & 50 & 100 & 500 & 1000 \\
\end{tabular}

A very common misunderstanding of the Roman numerals is that they use a universal subtractive syntax. In a subtractive syntax,
a lower digit written before a larger digit will be subtracted from the larger digit. 4 can be written as IIII and IV $(5-1)$. Romans 
themselves only really used the subtractive syntax for smaller numbers. For larger numbers, it was largely avoided, to give more clarity. 
For example, $499$ can be written LDVLIV, XDIX, VDIV or ID, while older sources most
likely would use CDXCIX. Ain't Roman numbers fun?

You are given $n$ numbers written in the Roman numeral system, and your task is to convert them
to the regular Arabic number system.

For the purposes of this problem we will subtract any digit written to the left of a larger digit, even if they are not directly adjacent.

